\section*{Bab \ref{ch:b1:arith} --- Aritmetika Bilangan Bulat}

\begin{solution}{exo:b1:arith:1}
$1001 = 1 \times 777 + 224$; $777 = 3 \times 224 + 105$; $224 = 2
\times 105 + 14$; $105 = 7 \times 14 + 7$; $14 = 2 \times 7 + 0$. Jadi
$\gcd(1001, 777) = 7$. Secara mundur:
\[
7 = 105 - 7 \times 14
= 105 - 7(224 - 2\times 105) = 15 \times 105 - 7 \times 224
\]
\[
= 15(777 - 3\times 224) - 7\times 224 = 15 \times 777 - 52 \times 224
= 15 \times 777 - 52(1001 - 777) = 67 \times 777 - 52 \times 1001 .
\]
Periksa: $67 \times 777 = 52\,059$ dan $52 \times 1001 = 52\,052$;
selisihnya $7$. Pasangan Bézout: $(u, v) = (-52, 67)$ untuk $1001u + 777v =
7$.
\end{solution}

\begin{solution}{exo:b1:arith:2}
Karena $10 \equiv 1 \pmod 9$: $10^k \equiv 1$, jadi $\sum_k d_k 10^k
\equiv \sum_k d_k \pmod 9$. Karena $10 \equiv -1 \pmod{11}$: $10^k
\equiv (-1)^k$, jadi bilangannya kongruen dengan jumlah berselang-seling
$\sum_k (-1)^k d_k$ modulo $11$ (dimulai dari angka \emph{satuannya}
dengan tanda $+$).

Untuk $123\,456\,789$: jumlah angkanya $45 \equiv 0 \pmod 9$. Jumlah berselang-selingnya dari
satuan: $9 - 8 + 7 - 6 + 5 - 4 + 3 - 2 + 1 = 5$, jadi bilangan itu
$\equiv 5 \pmod{11}$.
\end{solution}

\begin{solution}{exo:b1:arith:3}
Euclid: $237 = 2 \times 91 + 55$; $91 = 1 \times 55 + 36$; $55 = 1
\times 36 + 19$; $36 = 1 \times 19 + 17$; $19 = 1 \times 17 + 2$; $17
= 8 \times 2 + 1$. Secara mundur:
\[
1 = 17 - 8\times 2 = 17 - 8(19 - 17) = 9\times 17 - 8\times 19
= 9(36 - 19) - 8\times 19 = 9\times 36 - 17\times 19
\]
\[
= 9\times 36 - 17(55 - 36) = 26\times 36 - 17\times 55
= 26(91 - 55) - 17\times 55 = 26\times 91 - 43\times 55
\]
\[
= 26\times 91 - 43(237 - 2\times 91) = 112 \times 91 - 43 \times 237.
\]
Jadi $91 \times 112 \equiv 1 \pmod{237}$: penyelesaiannya adalah $x \equiv
112 \pmod{237}$. (Periksa: $91 \times 112 = 10\,192 = 43 \times 237 +
1$.)
\end{solution}

\begin{solution}{exo:b1:arith:4}
$\gcd(17, 39) = 1$: Euclid memberikan $39 = 2\times 17 + 5$, $17 = 3\times
5 + 2$, $5 = 2\times 2 + 1$, dan secara mundur
\[
1 = 5 - 2\times 2 = 5 - 2(17 - 3\times 5) = 7\times 5 - 2\times 17
= 7(39 - 2\times 17) - 2\times 17 = 7\times 39 - 16\times 17 .
\]
Penyelesaian khususnya $(x_0, y_0) = (-16, 7)$. Penyelesaian umum
persamaan homogennya $17x + 39y = 0$: $x = 39k$, $y = -17k$ (karena
$17 \mid 39y$ dan $\gcd(17,39) = 1$ memaksa $17 \mid y$ ---
menurut lema Gauss). Jadi
\[
(x, y) = (-16 + 39k,\; 7 - 17k), \qquad k \in \Z .
\]
Untuk ruas kanan $5$, kalikan penyelesaian khususnya dengan $5$:
$(x, y) = (-80 + 39k,\; 35 - 17k)$, $k \in \Z$.
\end{solution}

\begin{solution}{exo:b1:arith:5}
Untuk setiap \omterm{def:b1:arith:prime}{prima} $p$, dengan $\alpha = v_p(a)$ dan $\beta = v_p(b)$:
\[
v_p\bigl(\gcd(a,b)\bigr) + v_p\bigl(\operatorname{lcm}(a,b)\bigr)
= \min(\alpha, \beta) + \max(\alpha, \beta)
= \alpha + \beta = v_p(ab) .
\]
Dua bilangan bulat positif yang valuasinya sama pada setiap \omterm{def:b1:arith:prime}{prima} adalah sama
(\cref{prop:b1:arith:valuation}), jadi $\gcd(a,b)\operatorname{lcm}(a,b)
= ab$.
\end{solution}

\begin{solution}{exo:b1:arith:6}
Valuasinya: $\gcd(a, b) = 2^{\min(10,6)} 3^{\min(4,7)} 5^{\min(2,0)}
7^{\min(0,1)} = 2^6\, 3^4 = 5184$;
$\operatorname{lcm}(a,b) = 2^{10}\, 3^7\, 5^2\, 7$.

Cacah pembaginya: pembagi positif $n = \prod p_i^{\alpha_i}$ tepat
berupa pilihan $\prod p_i^{\beta_i}$ dengan $0 \leq \beta_i \leq
\alpha_i$ (\cref{prop:b1:arith:valuation}); pilihannya saling
bebas, jadi ada $\prod_i (\alpha_i + 1)$ pembagi. Untuk $a$:
$(10+1)(4+1)(2+1) = 165$.
\end{solution}

\begin{solution}{exo:b1:arith:7}
Andaikan $\sqrt p = \frac rq$ dengan $r, q \in \N^*$, yakni $p q^2 =
r^2$. Terapkan $v_p$: $v_p(pq^2) = 1 + 2v_p(q)$ ganjil, sedangkan $v_p(r^2)
= 2 v_p(r)$ genap. Sebuah bilangan bulat tak mungkin bervaluasi $p$ ganjil
sekaligus genap: kontradiksi. Jadi $\sqrt p \notin \Q$.
\end{solution}

\begin{solution}{exo:b1:arith:8}
\emph{Fakta umumnya.} Dengan $\gcd(m,n) = 1$, Bézout memberikan $mu + nv = 1$.
Tulis $x_0 = b\,mu + a\,nv$. Maka $x_0 \equiv a\,nv \equiv a(1 - mu)
\equiv a \pmod m$ dan serupa itu $x_0 \equiv b \pmod n$: itulah keberadaannya. Jika
$x$ dan $x'$ dua penyelesaian, maka $m$ dan $n$ \omterm{def:b1:arith:divides}{membagi} $x - x'$, sehingga
$mn \mid x - x'$ (\cref{thm:b1:arith:gauss}~(2)): itulah ketunggalannya modulo
$mn$.

\emph{Secara numerik:} $m = 7$, $n = 11$: $7 \times (-3) + 11 \times 2 =
1$. Jadi $x_0 = 5 \times 7 \times (-3) + 2 \times 11 \times 2 = -105 +
44 = -61 \equiv 16 \pmod{77}$. Periksa: $16 = 2\times 7 + 2 \equiv 2
\pmod 7$; $16 = 11 + 5 \equiv 5 \pmod{11}$. Penyelesaiannya: $x \equiv 16
\pmod{77}$.
\end{solution}

\begin{solution}{exo:b1:arith:9}
Modulo $7$: Fermat memberikan $3^6 \equiv 1$, dan $1000 = 6 \times 166 + 4$,
jadi $3^{1000} \equiv 3^4 = 81 \equiv 4 \pmod 7$.

Dua angka terakhir $7^{100}$: bekerjalah modulo $4$ dan modulo $25$. Modulo $4$: $7
\equiv -1$, jadi $7^{100} \equiv 1$. Modulo $25$: $7^2 = 49 \equiv -1$, jadi
$7^4 \equiv 1$ dan $7^{100} = (7^4)^{25} \equiv 1$. Menurut teorema sisa
Cina (\cref{exo:b1:arith:8}), $7^{100} \equiv 1
\pmod{100}$: jadi dua angka terakhirnya adalah $01$.
\end{solution}

\begin{solution}{exo:b1:arith:10}
Tulis $m = nq + r$, $0 \leq r < n$. Maka
\[
2^m - 1 = 2^r\bigl(2^{nq} - 1\bigr) + 2^r - 1,
\]
dan $2^n - 1$ \omterm{def:b1:arith:divides}{membagi} $2^{nq} - 1 = (2^n - 1)(2^{n(q-1)} + \dots +
1)$. Jadi modulo $2^n - 1$ berlaku $\;2^m - 1 \equiv 2^r - 1$, dan karena $0 \leq
2^r - 1 < 2^n - 1$, ini \emph{memang} sisa Euclidnya.

Karena itu \omterm{met:b1:arith:euclid}{algoritma Euclid} pada pasangan $(2^m - 1, 2^n - 1)$
mencerminkan, eksponen demi eksponen, algoritmanya pada $(m, n)$: setiap
langkah pembagiannya mengganti $(m, n)$ dengan $(n, r)$ di lantai atas dan $(2^m - 1,
2^n - 1)$ dengan $(2^n - 1, 2^r - 1)$ di lantai bawah. Algoritma di atasnya
berhenti pada $\gcd(m,n)$, jadi di bawahnya ia berhenti pada
$2^{\gcd(m,n)} - 1$.
\end{solution}

\begin{solution}{exo:b1:arith:11}
Selesaikan dulu $x^2 \equiv 1 \pmod p$: di sini $p \mid (x-1)(x+1)$, jadi menurut
lema Euclid $x \equiv 1$ atau $x \equiv -1 \pmod p$.

Pada hasil kali $(p-1)! = 1 \times 2 \times \dots \times (p-1)$, setiap
faktor $a$ terbalikkan modulo $p$, dan inversnya $a^{-1}$ kembali menjadi
salah satu faktornya (\cref{prop:b1:arith:invmod}). Pasangkan setiap $a$ dengan
$a^{-1}$: pasangannya berhasil kali $1$, kecuali faktor yang berpasangan dengan dirinya
sendiri ($a = a^{-1}$, yakni $a^2 \equiv 1$) yang berdiri sendirian --- dan itu
tepat $1$ dan $p - 1$. Jadi
\[
(p-1)! \equiv 1 \times (p - 1) \equiv -1 \pmod p .
\]
(Untuk $p = 2$: $1! = 1 \equiv -1 \pmod 2$; argumen pemasangannya
merosot tetapi hasilnya tetap berlaku.)

\emph{Konversnya.} Misalkan $n \geq 2$ komposit, $n = ab$ dengan $1 < a
\leq b < n$. Jika $a < b$, keduanya muncul sebagai faktor $(n-1)!$ yang berbeda,
sehingga $n \mid (n-1)!$ dan $(n-1)! \equiv 0 \not\equiv -1$. Jika $a = b$
(yakni $n = a^2$): untuk $a \geq 3$, baik $a$ maupun $2a$ bernilai $< n$, jadi
$n = a^2 \mid a \times 2a \mid (n-1)!$, dengan kesimpulan yang sama; adapun untuk $n = 4$,
$(n-1)! = 6 \equiv 2 \not\equiv -1 \pmod 4$.
\end{solution}

\begin{solution}{exo:b1:arith:12}
\begin{enumerate}
  \item Dengan induksi. Untuk $n = 1$: $F_0 = 3 = F_1 - 2 = 5 - 2$.
        Dengan mengandaikan $F_0\cdots F_{n-1} = F_n - 2$:
        \[
        F_0 \cdots F_n = (F_n - 2)F_n
        = \bigl(2^{2^n} - 1\bigr)\bigl(2^{2^n} + 1\bigr)
        = 2^{2^{n+1}} - 1 = F_{n+1} - 2 .
        \]
  \item Misalkan $m < n$ dan $d = \gcd(F_m, F_n)$. Menurut (1), $F_m$
        \omterm{def:b1:arith:divides}{membagi} $F_n - 2$, jadi $d$ \omterm{def:b1:arith:divides}{membagi} $F_n$ maupun $F_n -
        2$, sehingga ia \omterm{def:b1:arith:divides}{membagi} $2$. Tetapi setiap bilangan Fermat ganjil, jadi
        $d = 1$.
  \item Setiap $F_n \geq 3$ mempunyai pembagi \omterm{def:b1:arith:prime}{prima} $p_n$
        (menurut langkah pertama \cref{thm:b1:arith:euclidprimes}). Jika $m
        \neq n$, maka $p_m \neq p_n$, karena \omterm{def:b1:arith:prime}{prima} persekutuannya akan
        \omterm{def:b1:arith:divides}{membagi} $\gcd(F_m, F_n) = 1$. Jadi \omterm{def:b1:logic:map}{pemetaan} $n \mapsto p_n$
        bersifat \omterm{def:b1:logic:inj}{injektif} dari $\N$ ke dalam \omterm{def:b1:arith:prime}{bilangan prima}: sehingga \omterm{def:b1:arith:prime}{bilangan
        prima} ada tak hingga banyaknya.
\end{enumerate}
\end{solution}

\begin{solution}{pb:b1:arith:1}
\textbf{1.} $10! = 3\,628\,800$: dengan dua nol penutup. Valuasinya
faktor demi faktor: pangkat $2$ datang dari $2, 4 = 2^2, 6, 8 = 2^3,
10$, sehingga totalnya $v_2(10!) = 1 + 2 + 1 + 3 + 1 = 8$; sedangkan pangkat $5$
datang dari $5$ dan $10$: $v_5(10!) = 2$. Nol penutupnya $=
\min(v_2, v_5) = 2$, yang sejalan.

\textbf{2.} Tulis pembagian Euclid $\lfloor x\rfloor = nq +
r$, $0 \leq r \leq n - 1$. Maka $x = nq + r + \{x\}$ dengan $0 \leq
r + \{x\} < n$, sehingga $\lfloor x/n \rfloor = q = \bigl\lfloor \lfloor
x \rfloor / n \bigr\rfloor$.

\textbf{3.} Misalkan $\alpha = v_p(a) \leq \beta = v_p(b)$ (tukarkan bila
perlu) lalu tulis $a = p^\alpha a'$, $b = p^\beta b'$ dengan $p
\nmid a', b'$. Maka $a + b = p^\alpha\bigl(a' + p^{\beta -
\alpha}b'\bigr)$, jadi $v_p(a + b) \geq \alpha = \min$. Jika $\alpha <
\beta$, kurungnya bernilai $a' + p^{\beta-\alpha}b' \equiv a'
\not\equiv 0 \pmod p$: sehingga valuasinya tepat $\alpha$.

\textbf{4.} Kelipatan $m$ di $\intint1n$ adalah $m, 2m, \dots,
qm$ dengan $q$ bilangan bulat terbesar yang memenuhi $qm \leq n$, yakni $q =
\lfloor n/m \rfloor$.

\textbf{5.} Menurut ketunggalan pemfaktoran, $v_p(n!) = \sum_{j=1}^{n}
v_p(j)$. Cacahlah secara berbeda: setiap $j$ menyumbang $v_p(j) =
\#\{k \geq 1 : p^k \mid j\}$, sehingga
\[
v_p(n!) = \sum_{j=1}^n \#\{k : p^k \mid j\}
= \sum_{k\geq1} \#\{j \leq n : p^k \mid j\}
= \sum_{k\geq1} \Bigl\lfloor \frac n{p^k} \Bigr\rfloor
\]
menurut pertanyaan 4 --- itulah rumus Legendre. Jumlahnya hingga: suku
dengan $p^k > n$ lenyap.

\textbf{6.} $v_5(1000!) = 200 + 40 + 8 + 1 = 249$ (dari pembagian dengan
$5, 25, 125, 625$); $v_2(1000!) = 500 + 250 + 125 + 62 + 31 + 15 +
7 + 3 + 1 = 994$. Nol penutup $1000!$: setiap nol menghabiskan
satu faktor $2$ dan satu faktor $5$, jadi ada $\min(994, 249) = 249$ nol.

\textbf{7.} Dengan $n = \sum_i a_ip^i$, pertanyaan 2 memberikan $\lfloor
n/p^k \rfloor = \sum_{i \geq k} a_ip^{i-k}$ (yakni memenggal ekspansi
basis $p$). Dengan menjumlahkan atas $k \geq 1$ lalu menukarkan kedua
jumlah hingga itu:
\[
v_p(n!) = \sum_{i\geq1} a_i \sum_{k=1}^{i} p^{i-k}
= \sum_{i\geq0} a_i\,\frac{p^i - 1}{p - 1}
= \frac{n - s_p(n)}{p - 1} .
\]

\textbf{8.} Untuk $p = 2$: $v_2(n!) = n - s_2(n)$. Karena $n \geq 1$
mempunyai $s_2(n) \geq 1$, selalu berlaku $v_2(n!) \leq n - 1 < n$: jadi $2^n \nmid
n!$. Dan $v_2(n!) = n - 1$ jika dan hanya jika $s_2(n) = 1$, yakni jika dan hanya jika $n$ pangkat
$2$.

\textbf{9.} Di sini $n$ mempunyai $\lfloor \log_p n \rfloor + 1$ angka
basis $p$, masing-masing paling besar $p - 1$, jadi $1 \leq s_p(n) \leq
(p-1)\bigl(\log_p(n) + 1\bigr)$. Dengan mensubstitusikannya pada pertanyaan 7:
\[
\frac n{p-1} - \log_p(n) - 1 \;\leq\; v_p(n!) \;<\; \frac n{p-1},
\]
lalu dengan membaginya dengan $n$: $\frac{v_p(n!)}n \to \frac1{p-1}$.

\textbf{10.} $Z(n) - Z(n-1) = v_5(n!/(n-1)!) = v_5(n)$: jadi cacah
nol penutupnya melonjak sebesar $v_5(n)$ pada setiap kelipatan $5$ dan
tetap di antaranya. Di sini $Z(24) = \lfloor24/5\rfloor = 4$ dan $Z(25) =
5 + 1 = 6$: pada $n = 25$ cacahnya melonjak dari $4$ langsung ke $6$
(karena $v_5(25) = 2$), dan karena $Z$ tidak turun dengan $Z \leq 4$
sebelumnya dan $Z \geq 6$ sesudahnya, nilai $5$ tak pernah tercapai: jadi tak ada
faktorial yang berakhir dengan tepat lima nol.

\textbf{11.} Tulis $x = \lfloor x\rfloor + \{x\}$: maka $\lfloor x +
y\rfloor = \lfloor x\rfloor + \lfloor y\rfloor + \lfloor \{x\} +
\{y\}\rfloor$, dan $0 \leq \{x\} + \{y\} < 2$ membuat lantai terakhirnya
$0$ atau $1$. Lalu, dengan Legendre yang diterapkan tiga kali,
\[
v_p\binom{m+n}m = v_p\bigl((m{+}n)!\bigr) - v_p(m!) - v_p(n!)
= \sum_{k\geq1}\Bigl(
\Bigl\lfloor\frac{m+n}{p^k}\Bigr\rfloor
- \Bigl\lfloor\frac{m}{p^k}\Bigr\rfloor
- \Bigl\lfloor\frac{n}{p^k}\Bigr\rfloor\Bigr),
\]
yaitu jumlah hingga berisi $0$ dan $1$ (terapkan klaim pertamanya pada $x =
m/p^k$, $y = n/p^k$).

\textbf{12.} Tetapkan $k \geq 1$ lalu tulis $m = p^km_1 + m_0$, $n =
p^kn_1 + n_0$ dengan $0 \leq m_0, n_0 < p^k$ (lewat pembagian Euclid:
$m_0$ adalah bilangan yang dibentuk oleh $k$ angka rendah $m$). Maka
\[
\Bigl\lfloor\frac{m+n}{p^k}\Bigr\rfloor
- \Bigl\lfloor\frac m{p^k}\Bigr\rfloor
- \Bigl\lfloor\frac n{p^k}\Bigr\rfloor
= \Bigl\lfloor\frac{m_0 + n_0}{p^k}\Bigr\rfloor ,
\]
yang bernilai $1$ bila $m_0 + n_0 \geq p^k$ dan $0$ bila tidak. Tetapi $m_0 +
n_0 \geq p^k$ mengatakan persis bahwa menjumlahkan $k$ angka rendah $m$
dan $n$ meluap ke posisi $k$ --- yaitu sebuah simpanan ke posisi $k$
pada algoritma penjumlahan seperti di sekolah. Dengan menjumlahkan atas $k$:
jadi $v_p\binom{m+n}m$ adalah banyaknya simpanan pada penjumlahan basis $p$
untuk $m + n$. (Kummer, 1852.)

\textbf{13.} Terapkan Kummer pada $m = j$, $n = p^k - j$, dengan jumlah $p^k =
(1\underbrace{0\cdots0}_{k})_p$. Misalkan $a = v_p(j)$, sehingga angka
basis $p$ dari $j$ pada posisi $0, \dots, a-1$ semuanya $0$ dan angka
pada posisi $a$ tak nol. Angka $p^k - j$ di bawah posisi
$a$ juga $0$ (karena $p^k - j = p^a(p^{k-a} - j/p^a)$). Pada posisi
$a$, kedua angka tak nolnya harus berjumlah $p$ (agar angka hasilnya $0$):
yaitu satu simpanan; lalu pada setiap posisi $a+1, \dots, k-1$, angkanya beserta
simpanan yang masuk berjumlah $p$ (angka hasilnya $0$ lagi): jadi simpanannya
merambat. Totalnya: $k - a$ simpanan, sehingga $v_p\binom{p^k}j = k -
v_p(j)$. Untuk $k = 1$: $v_p\binom pj = 1$ untuk $0 < j < p$, yaitu
\omterm{def:b1:arith:divides}{keterbagian} yang dipakai pada \cref{thm:b1:arith:fermat}.

\textbf{14.} Menurut bentuk digitalnya (pertanyaan 7), dengan memakai $s_2(2n) =
s_2(n)$ (karena menambahkan satu angka nol):
\[
v_2\binom{2n}n = \bigl(2n - s_2(2n)\bigr) - 2\bigl(n -
s_2(n)\bigr) = 2s_2(n) - s_2(2n) = s_2(n) \geq 1 :
\]
$\binom{2n}n$ selalu genap, dan $v_2 = 1$ (yakni $\binom{2n}n
\equiv 2 \pmod 4$) tepat ketika $s_2(n) = 1$, yakni ketika $n$ merupakan
pangkat $2$.

\textbf{15.} Vandermonde dengan $m = n = k = p$: $\binom{2p}p =
\sum_{j=0}^p \binom pj\binom p{p-j} = \sum_{j=0}^p \binom pj^2$.
Untuk $0 < j < p$ berlaku $p \mid \binom pj$ (pertanyaan 13), jadi $\binom
pj^2 \equiv 0 \pmod p$; sedangkan suku ujungnya memberikan $1 + 1$:
sehingga $\binom{2p}p \equiv 2 \pmod p$.

\textbf{16.} Basis $3$: $500 = 486 + 9 + 3 + 2$, dengan angka (dari rendah ke
tinggi) $(2, 1, 1, 0, 0, 2)$, jadi $s_3(500) = 6$; dan $1000 = 729 +
243 + 27 + 1$, dengan angka $(1, 0, 0, 1, 0, 1, 1)$, jadi $s_3(1000) = 4$.
\emph{Kummer:} jumlahkan $500 + 500$ dalam basis $3$: posisi $0$: $2 + 2 =
4$, angkanya $1$ dengan simpanan $1$; posisi $1$: $1 + 1 + 1 = 3$, angkanya $0$
dengan simpanan $1$; posisi $2$: $1 + 1 + 1 = 3$, angkanya $0$ dengan simpanan $1$;
posisi $3$: $0 + 0 + 1 = 1$, tanpa simpanan; posisi $4$: $0$;
posisi $5$: $2 + 2 = 4$, angkanya $1$ dengan simpanan $1$; posisi $6$:
simpanannya mendarat: angkanya $1$. Empat simpanan: jadi $v_3\binom{1000}{500} = 4$.
\emph{Legendre:} $v_3(1000!) = \frac{1000 - 4}2 = 498$ dan
$v_3(500!) = \frac{500 - 6}2 = 247$, sehingga $v_3\binom{1000}{500} =
498 - 2\times247 = 4$. Kedua perhitungannya cocok --- dan
angka penjumlahannya $(1, 0, 0, 1, 0, 1, 1)$ menghasilkan $1000$, sebagaimana
seharusnya.

\textbf{17.} Menurut Kummer (dengan $p = 2$, $m = k$, $n' = n - k$):
$\binom nk$ ganjil jika dan hanya jika penjumlahan $k + (n - k)$ dalam basis $2$ tak
mempunyai simpanan, yakni jika dan hanya jika pada setiap posisi angkanya memenuhi $k_i + (n -
k)_i = n_i$; dan dalam hal itu $k_i \leq n_i$ untuk setiap $i$. Sebaliknya,
jika $k_i \leq n_i$ untuk setiap $i$, maka bilangan berangka $n_i -
k_i$ sama dengan $n - k$ dan penjumlahannya bebas simpanan. Bukti yang sama dalam basis
$p$: $p \nmid \binom nk$ jika dan hanya jika setiap angka basis $p$ dari $k$ paling
besar sama dengan angka $n$ yang bersesuaian.

\textbf{18.} Dengan mencacah $k \in \intint0n$ yang angkanya menuruti
$k_i \leq n_i$: setiap angka $k$ dipilih secara bebas di antara
$n_i + 1$ nilai, sehingga ada $\prod_i (n_i + 1)$ pilihan; dan dalam basis $2$
ini sama dengan $2^{\#\{i : n_i = 1\}} = 2^{s_2(n)}$. Baris $4 =
(100)_2$: ada $2^1 = 2$ entri ganjil --- memang $1, 4, 6, 4, 1$ berentri
ganjil hanya di kedua ujungnya. Baris $5 = (101)_2$: ada $2^2 = 4$ --- memang
$1, 5, 10, 10, 5, 1$.

\textbf{19.} Semua entri bagian dalamnya genap $\iff$ barisnya mempunyai tepat
$2$ entri ganjil (karena kedua ujungnya selalu ganjil) $\iff 2^{s_2(n)} =
2 \iff s_2(n) = 1 \iff n$ merupakan pangkat $2$.

\textbf{20.} Pada jumlah pertanyaan 11, suku ke-$k$ lenyap begitu
$p^k > m + n$ (karena ketiga lantainya lalu sama; memang yang
pertama $0$ bila $p^k > m+n$; lebih sederhananya setiap sukunya $0$).
Jadi paling banyak $\lfloor \log_p(m+n)\rfloor$ suku yang tak nol, masing-masing
bernilai $1$: sehingga $a = v_p\binom{m+n}m \leq \log_p(m+n)$, yakni $p^a \leq
m + n$.

\textbf{21.} Untuk setiap \omterm{def:b1:arith:prime}{prima} $p$ berlaku $v_p\bigl(\operatorname{lcm}(1,
\dots, 2n)\bigr) = \lfloor\log_p(2n)\rfloor$ (karena pangkat terbesar
$p$ yang tak melampaui $2n$ muncul di antara $1, \dots, 2n$). Pertanyaan 20
dengan $m = n$ memberikan $v_p\binom{2n}n \leq \lfloor\log_p(2n)\rfloor$
untuk setiap $p$: jadi menurut \cref{prop:b1:arith:valuation}, $\binom{2n}n
\mid \operatorname{lcm}(1, \dots, 2n)$. Untuk ukurannya: rasio
$\binom{2n}{k+1}/\binom{2n}k = \frac{2n-k}{k+1} \geq 1$ tepat
untuk $k < n$, jadi entri pusatnya adalah yang terbesar di antara $2n + 1$
entri baris $2n$, sehingga $4^n = \sum_k \binom{2n}k \leq
(2n+1)\binom{2n}n$. Dengan menggabungkannya:
\[
\operatorname{lcm}(1, \dots, 2n) \geq \binom{2n}n \geq
\frac{4^n}{2n + 1} .
\]
Seandainya \omterm{def:b1:arith:prime}{bilangan prima} di bawah $2n$ hanya sedikit, KPK-nya tak akan sebesar
itu: jadi pertumbuhan eksponensial KPK-nya adalah jejak kuantitatif dari
limpahnya \omterm{def:b1:arith:prime}{bilangan prima}.

\textbf{22.} $Z(n) = \sum_k\lfloor n/5^k\rfloor \approx \frac
n4$, jadi bidiklah di dekat $n = 4 \times 2026 = 8104$: $Z(8104) = 1620 +
324 + 64 + 12 + 2 = 2022$. Naikkan dengan kelipatan $5$: $Z(8110)
= 2024$, $Z(8115) = 2025$, dan
\[
Z(8120) = 1624 + 324 + 64 + 12 + 2 = 2026 .
\]
Karena $Z$ tetap di antara kelipatan $5$ dan $Z(8119) =
Z(8115) = 2025$, maka $n$ terkecil yang berisi sekurang-kurangnya $2026$ nol
penutup adalah $n = 8120$.

\textbf{23.} Basis $7$: $50 = 49 + 1$, dengan angka (dari rendah ke tinggi) $(1, 0,
1)$. Menjumlahkan $50 + 50$: posisi $0$: $1 + 1 = 2 < 7$, tanpa simpanan;
posisi $1$: $0 + 0 = 0$; posisi $2$: $1 + 1 = 2 < 7$, tanpa
simpanan. Bebas simpanan, jadi menurut Kummer $v_7\binom{100}{50} = 0$: sehingga $7 \nmid
\binom{100}{50}$. Legendre sependapat: $v_7(100!) = \lfloor 100/7
\rfloor + \lfloor 100/49 \rfloor = 14 + 2 = 16$ dan $v_7(50!) = 7
+ 1 = 8$, jadi $v_7\binom{100}{50} = 16 - 2\times8 = 0$.

\textbf{24.} (i) Ketunggalan pemfaktoran melandasi definisi
$v_p$ itu sendiri beserta keaditifannya, sehingga juga rumus Legendre
dan setiap kesimpulan \omterm{def:b1:arith:divides}{keterbagiannya}
(\cref{prop:b1:arith:valuation}). (ii) Pembagian Euclid menghasilkan
kesamaan pemenggalan pada pertanyaan 2 dan pemilahan $m = p^km_1 +
m_0$ yang mengisolasi simpanannya (pertanyaan 12). (iii) Pencacahan: cacah
kelipatan $m$ (pertanyaan 4), hasil kali pilihan angkanya
(pertanyaan 18), dan batas jumlah barisnya $4^n \leq
(2n+1)\binom{2n}n$ (pertanyaan 21) semuanya argumen
bergaya \cref{ch:b1:counting}.

\textbf{25.} Legendre mengubah ``pangkat $p$ yang mana yang \omterm{def:b1:arith:divides}{membagi}
$n!$'' menjadi aritmetika angka basis $p$; Kummer memampatkan
jawabannya untuk \omterm{def:b1:counting:objects}{koefisien binomial} menjadi simpanan pada satu
penjumlahan --- \omterm{def:b1:arith:divides}{keterbagian}, yang tampaknya sifat global bilangan raksasa,
ternyata terbaca secara lokal, angka demi angka. Pertanyaan 16 adalah
paradigmanya: empat simpanan, yang dihitung dengan tangan, menentukan pangkat
$3$ yang persis pada bilangan beratus angka. Dan pertanyaan 21
memperlihatkan lingkaran gagasan yang sama menyenggol perairan dalam: batas
bawah eksponensial bagi $\operatorname{lcm}(1, \dots, 2n)$ adalah
langkah pertama yang sepenuhnya dasar menuju teorema \omterm{def:b1:arith:prime}{bilangan prima},
yang buktinya jauh di luar jilid ini. Seluruh kotak perkakasnya ---
pembagian, FPB, valuasi --- diputar ulang untuk polinomial pada
\cref{ch:b1:poly}, yang di sana analog ekspansi angkanya adalah
ekspansi dalam pangkat $(X - a)$.

\end{solution}
